\section{Artifact and comparison details}
\label{app:evidence}
The companion \texttt{ARTIFACT.md} gives executable commands, source pins,
dependency locks and the mapping from protocol objects to implementation.
The claim register records assumptions and counterexamples for each claim.
Current native evidence and historical experiments retain separate source
inventories. The PDF build verifies the derived results and frozen artifacts.

\subsection{Current native qualification}
The current campaign records 21 terminal commands, 36,556 unchanged source
inputs and 48 hashed outputs. Rust pass counts by boundary are 54 workflow,
10 native delegation labels, five existing native composition, 453
swarm/runtime, 11 execution-export, 25 durable SQLite and 91 default standalone
cases. The delegation count includes its subprocess worker.
One pre-existing workflow doctest is ignored. The default standalone run
ignores six opt-in chain cases; a separate run passes all six against the
identical source inventory, exercising 97 standalone cases in total.

The campaign includes three Node wire/pin tests, 13 artifact tests, the evolving
funded execution, four existing earned-child crash subcases, four strict
Clippy commands, three format commands, the delegation example and the funded
binary build. This qualifies the affected local boundaries, not the full
workspace or a production deployment.

The suite outputs identify the tested bindings and recovery cuts. Intended
negative controls and initial development failures remain distinct from the
terminal passing results.

\subsection{Retained financial model and implementation checks}
\begin{table}[ht]
\centering\small
\begin{tabularx}{\linewidth}{@{}l r X@{}}
\toprule
Experiment & Result & Scope\\
\midrule
Finite claim graph & \ModelStates\ states & Two prefunded edges and nine timestamps\\
Transition checks & \ModelTransitions & Selected role, deadline and financial transitions\\
Contract suite & \ContractTests\ passing & Private-chain bytecode; \TraceCount\ representative traces\\
SQL comparator & \MatchedSteps\ matched steps & Same representative traces; trusted sequential store\\
Upstream ERC-8183 & \UpstreamTests\ passing & Unmodified suite at the pinned source\\
Paired escrow cases & \PairedTests\ passing & Same asset, fees, capital and evaluator trust\\
Lean specification & \ProofCount\ results & Abstract conservation and earned-state lemmas\\
\bottomrule
\end{tabularx}
\caption{Historical results retain their original source identities. Their
counts are not attributed to the changed native implementation.}
\label{tab:evidence}
\end{table}

The finite graph begins with a three-unit parent and a two-unit child.
It explores submissions, decisions, payment, refund, expiry and execution
uncertainty within that alphabet. The Solidity runner replays \TraceCount\
representative financial traces, rather than all \ModelTransitions\
transitions. It compares state, commitments, decisions and actual balances.
A deliberately weakened model permits an earned claim to refund; the explorer
finds the counterexample. An in-memory mutation of the escrow produces the
corresponding intended test failure.

The Lean specification proves conservation over finite account traces,
paid-plus-refunded bounds, persistence of earned states, and noninterference
between distinct allocation entries. The native binding and graph-growth
arguments are not mechanized refinements of the Rust code. These different
forms of evidence support their stated boundaries rather than a proof of the
implementation as a whole.

\subsection{Matched alternatives}
The ordinary SQL comparator implements the same monetary transitions without
Chio kernel code. It matches \MatchedSteps\ non-time steps over the
\TraceCount\ representative traces, with additional tests for child collection,
allocation exclusivity and unknown execution. A conditional-backing experiment
also compares a structural bound with an independent exhaustive oracle over
4,095 finite profiles. A conventional allocator obtains the same capital
saving from selecting among exclusive alternatives before dispatch.

Upstream ERC-8183 is pinned to commit \texttt{142e669c1f}. Its
\UpstreamTests\ unchanged tests and \PairedTests\ added comparison tests use
equal mock assets, zero fees, capital and evaluator trust. Both implementations
pay a separately funded child after parent refund, reject two 60-unit jobs
funded from 100, and deny spending an unsettled parent reserve.
The upstream implementation also pays an approved 60-unit partial claim from a
100-unit job and refunds the remaining 40 after rejection.

The lifecycle differences are concrete. Chio records acceptance before later
withdrawal, applies a resolution cutoff, and preserves earned withdrawal through
its funding pause or allowlist changes. The pinned upstream completion pays
immediately; after its expiry grace period, completion can still win if refund
has not executed. Its refunds are blocked while paused. These observations
concern the pinned implementations, whose specifications and extension
mechanisms admit other choices. No performance or security ranking is inferred.

The historical security-integration record binds its native runs to their
earlier v2/v3 source inventories. Current composed qualification remains separate.
